\documentclass[a4paper,10.5pt]{article}
\usepackage[utf8]{inputenc}
\usepackage[T1]{fontenc}
\usepackage[french]{babel}
\usepackage[left=1cm,right=1cm,top=.5cm,bottom=0cm]{geometry}
\usepackage{enumitem}
\usepackage{hyperref}
\usepackage{xcolor}
\usepackage{titlesec}
\usepackage{fontawesome5}
\usepackage{lmodern}
\usepackage[normalem]{ulem}

% --- Couleurs Framatome ---
\definecolor{primary}{RGB}{200, 50, 30}
\definecolor{secondary}{RGB}{80, 80, 80}

% --- Config Liens ---
\hypersetup{colorlinks=true, linkcolor=primary, filecolor=primary, urlcolor=primary}

% --- Titres ---
\titleformat{\section}{\large\bfseries\color{primary}\uppercase}{}{0em}{}[\titlerule]
\titlespacing{\section}{0pt}{8pt}{5pt}

% --- Commandes ---
\newcommand{\resumeItem}[1]{\item\small{#1}}

\begin{document}

% --- EN-TÊTE ---
\begin{center}
    {\Large \textbf{Loïc Aron MBASSI EWOLO}} \\ \vspace{4pt}
    \textbf{\large Alternance Développeur Logiciel Embarqué} \\
    \textit{\small Disponible dès septembre 2026 pour 12 mois} \\ \vspace{4pt}
    \small
    \faMapMarker\ Cergy, France (Mobile France) \quad
    \faPhone\ (+33) 7 59 52 33 98 \quad
    {\color{primary}\faEnvelope\ \href{mailto:loicaron.mbassiewolo@ensea.fr}{loicaron.mbassiewolo@ensea.fr}} \\ \vspace{2pt}
    {\color{primary}\faLinkedin\ \href{https://www.linkedin.com/in/loic-mbassi/}{linkedin.com/in/loic-mbassi}} \quad
    {\color{primary}\faGithub\ \href{https://github.com/Nameless0l}{github.com/Nameless0l}} \quad
    {\color{primary}\faGlobe\ \href{https://mbassiloic.tech}{mbassiloic.tech}}
\end{center}

\vspace{-6pt}

% --- PROFIL ---
\noindent\small{Détection et isolation de défauts (FDI) sur drone avec contrôle tolérant aux pannes (FTC) en environnement safety-critical (MATLAB/Simulink). Développement embarqué C/C++ sur Nvidia Jetson (SLAM, Lidar 3D, YOLOv10) et prototypage IoT sur STM32 (capteurs IMU, fusion données). Double diplôme ENSEA/Polytechnique Yaoundé en systèmes embarqués, signal et IA. Programmation système C/Linux (IPC, mémoire partagée, signaux POSIX). Rigueur de codage, analyse statique et tests unitaires.}

% --- FORMATION ---
\section{Formation}
\begin{itemize}[leftmargin=*, label=\tiny$\bullet$, nosep]
    \item \textbf{ENSEA -- École Nationale Supérieure de l'Électronique et de ses Applications} \hfill \textit{2025 -- 2027} \\
    \small{Double diplôme d'Ingénieur -- Systèmes Embarqués, FPGA/VHDL, Traitement du Signal, Deep Learning.}
    \item \textbf{École Polytechnique de Yaoundé (1\textsuperscript{ère} école d'ingénieur CEMAC)} \hfill \textit{2021 -- 2025} \\
    \small{Diplôme d'Ingénieur de Conception (Master 2) : \textbf{Mathématiques Appliquées}, Génie Logiciel, Algorithmique.}
\end{itemize}

% --- COMPÉTENCES ---
\section{Compétences Techniques}
\begin{itemize}[leftmargin=*, nosep]
    \item \small{\textbf{Embarqué \& Sûreté :} C embarqué, C++, STM32, Arduino, Raspberry Pi, Nvidia Jetson, FPGA/VHDL (Quartus, ModelSim).}
    \item \small{\textbf{Tests \& Qualité logicielle :} Analyse statique de code, revues de code, tests unitaires, tests d'intégration, FDI/FTC.}
    \item \small{\textbf{OS \& Temps Réel :} Linux (IPC, SHM, signaux POSIX), ordonnancement FIFO, contraintes temps réel.}
    \item \small{\textbf{IHM \& Outils :} Python (Tkinter, scripting), Qt/C++ (notions), MATLAB/Simulink, Docker, Git.}
    \item \small{\textbf{Documentation :} Rédaction de spécifications techniques, documentation scientifique (français/anglais).}
    \item \small{\textbf{Langues :} Français (natif), Anglais (professionnel/technique).}
\end{itemize}

% --- EXPÉRIENCES / PROJETS ---
\section{Expériences \& Projets Embarqués}
\begin{itemize}[leftmargin=*, label={-}]

\item \textbf{Fault Detection \& Fault-Tolerant Control -- Drone (Safety-Critical)} \hfill \textit{2025} \\
\textit{Projet d'option -- Contrôle automatique, ENSEA} \hfill \textit{MATLAB/Simulink}
\vspace{-2pt}
    \begin{itemize}[leftmargin=0.4cm, nosep, label=\tiny$\bullet$]
        \item \small{Modélisation d'un système embarqué en conditions dégradées : perte d'efficacité d'actionneurs, biais et bruit capteurs.}
        \item \small{Conception d'un schéma FDI basé sur des observateurs (résidus, équations de parité), analyse de sensibilité.}
        \item \small{Stratégie FTC : reconfiguration de contrôleur, gain scheduling, modes de contrôle dégradés.}
        \item \small{Simulation de scénarios de pannes réalistes, évaluation comparative des performances.}
    \end{itemize}
\vspace{3pt}

\item \textbf{Challenge CoHoMa -- Robotique Autonome (ENSEA / Armée française)} \hfill \textit{2025 -- En cours} \\
\textit{Développement embarqué C/C++} \hfill \textit{C/C++, Nvidia Jetson, YOLOv10, SLAM, Lidar 3D, Docker}
\vspace{-2pt}
    \begin{itemize}[leftmargin=0.4cm, nosep, label=\tiny$\bullet$]
        \item \small{Développement C/C++ sur Jetson : cartographie autonome (SLAM), fusion capteurs Lidar VLP16/caméra de profondeur.}
        \item \small{Optimisation d'inférence IA embarquée sur GPU, containerisation Docker pour reproductibilité multi-plateforme.}
    \end{itemize}
\vspace{3pt}

\item \textbf{Simulation Flotte de Taxis -- Programmation Système C/Linux} \hfill \textit{2024} \\
\textit{Projet académique -- ENSPY} \hfill \textit{C, IPC (SHM, Signaux), OpenGL, Linux}
\vspace{-2pt}
    \begin{itemize}[leftmargin=0.4cm, nosep, label=\tiny$\bullet$]
        \item \small{Simulation multiprocessus en C : mémoire partagée, signaux POSIX, ordonnanceur FIFO.}
        \item \small{Gestion manuelle de la mémoire, synchronisation de processus concurrents, rigueur de programmation bas niveau.}
    \end{itemize}
\vspace{3pt}

\item \textbf{Harmony Gloves -- Gants IoT (STM32, Capteurs)} \hfill \textit{2024} \\
\textit{Labo IOTA, École Polytechnique de Yaoundé} \hfill \textit{STM32, Arduino, C, Capteurs IMU/Flex}
\vspace{-2pt}
    \begin{itemize}[leftmargin=0.4cm, nosep, label=\tiny$\bullet$]
        \item \small{Participation à la conception de gants instrumentés : lecture capteurs, fusion données, soudure composants.}
        \item \small{Intégration hardware/software, modèle ML embarqué pour classification gestuelle temps réel.}
    \end{itemize}
\vspace{3pt}

\item \textbf{Stage Recherche IoT -- Edge Computing \& TinyML} \hfill \textit{Oct 2023 -- Jan 2024} \\
\textit{IoT Lab, École Polytechnique de Yaoundé} \hfill \textit{Python, TF Lite, Arduino, Raspberry Pi, C}
\vspace{-2pt}
    \begin{itemize}[leftmargin=0.4cm, nosep, label=\tiny$\bullet$]
        \item \small{Optimisation de modèles ML pour microcontrôleurs : quantification, pruning, contraintes mémoire et temps réel.}
    \end{itemize}
\vspace{3pt}

\end{itemize}

% --- DISTINCTIONS ---
\section{Distinctions \& Compléments}
\begin{itemize}[leftmargin=*, label=\tiny$\bullet$, nosep]
    \item \small{\textbf{Stage INRIA Saclay} (Équipe TRIBE, avr--août 2026) : pipeline Python, analyse statique de code, rédaction scientifique anglais.}
    \item \small{\textbf{1\textsuperscript{er} Prix} CITS National 2024 (75 équipes). \textbf{Lead AI Cell} (ENSPY) : workshops embarqué, ML.}
    \item \small{\textbf{Certif.} Supervised ML (Stanford/DeepLearning.AI), Python for Data Science (IBM/Coursera).}
\end{itemize}

\end{document}
